\textbf{Proof.}
Fix \(v\in\mathcal H^s(2M)\) with \(\|v\|_\infty\le1\), and put \(u_{\ell,r}=v^{(r)}(z_\ell)\). The box restrictions follow immediately from the defining derivative bounds. For \(r=p\), the H\"older condition gives
\[
|u_{\ell,p}-u_{k,p}|\le2M|z_\ell-z_k|^\beta
=2M\frac{\Gamma(\beta+1)}{\Gamma(s-p+1)}|z_\ell-z_k|^{s-p}.
\tag{1}
\]
For \(0\le r<p\), repeated use of the fundamental theorem of calculus gives the Taylor formula with integral remainder
\[
v^{(r)}(x)-\sum_{q=r}^p v^{(q)}(y)\frac{(x-y)^{q-r}}{(q-r)!}
=\frac1{(p-r-1)!}\int_y^x(x-z)^{p-r-1}\{v^{(p)}(z)-v^{(p)}(y)\}\,dz,
\tag{2}
\]
with the oriented-integral interpretation when \(x<y\). Bounding the integrand by \(2M|z-y|^\beta\) and evaluating the beta integral yields
\[
\left|v^{(r)}(x)-\sum_{q=r}^p v^{(q)}(y)\frac{(x-y)^{q-r}}{(q-r)!}\right|
\le2M\frac{\Gamma(\beta+1)}{\Gamma(p-r+\beta+1)}|x-y|^{p-r+\beta}.
\tag{3}
\]
Because \(p+\beta=s\), (3) is precisely the jet constraint in the definition. Equations (1)--(3) prove feasibility for every admissible trace.

For such a trace,
\[
\sum_{i=1}^n\Gamma_i v(X_i)-\frac1m\sum_{j=1}^mv(X_j^T)
=\sum_{\ell=1}^Nc_\ell u_{\ell,0}.
\tag{4}
\]
Taking the supremum of the two signs in (4) over admissible functions, and then enlarging from admissible traces to the jet polytope, proves \(B_X\le\overline B_X\). There are \(N(p+1)\) variables. Splitting each absolute-value restriction into two one-sided inequalities gives \(O\{N^2(p+1)\}\) inequalities, and the box constraints make both objective values finite.

It remains to verify the exact low-smoothness statement. If \(0<s\le1\), then \(p=0\), \(\beta=s\), and feasibility is exactly
\[
|u_{\ell,0}|\le1,
\qquad
|u_{\ell,0}-u_{k,0}|\le2M|z_\ell-z_k|^s.
\tag{5}
\]
The function \(d_s(x,y)=|x-y|^s\) is a metric on \([0,1]\). The McShane formula
\[
\widetilde v(x)=\inf_{1\le\ell\le N}\{u_{\ell,0}+2M d_s(x,z_\ell)\}
\tag{6}
\]
extends the finite trace and has H\"older seminorm at most \(2M\). Projection of \(\widetilde v\) onto \([-1,1]\) preserves both its values at the sites and its H\"older constant. Thus every feasible jet in (5) is the trace of an admissible function. The relaxation is therefore exact for \(0<s\le1\). No extension assertion is made for the higher-order jet polytope. \(\square\)